%% Rendement global d'une chaîne d'énergie : les rendements se multiplient.
%% La bande orange (type Sankey) s'amincit à chaque bloc : la puissance transmise
%% diminue, la différence part en pertes.
\documentclass[tikz,border=4pt]{standalone}
\usepackage{liaisons}
\begin{document}
\begin{tikzpicture}[x=1cm,y=1cm,
  bloc/.style={draw=solideD, line width=1.1pt, fill=white, rounded corners=2pt,
               minimum width=1.75cm, minimum height=0.95cm, inner sep=1pt,
               align=center, font=\small},
  perte/.style={-{Stealth[length=5pt,width=4pt]}, line width=0.9pt, draw=black!55},
  et/.style={font=\scriptsize, inner sep=2pt}]

%% ================== bande de puissance (Sankey) =======================
\fill[solideD!35] (0.75,2.78) -- (7.55,2.30) -- (7.55,1.70) -- (0.75,1.22) -- cycle;
\fill[solideD!35] (7.55,2.50) -- (8.30,2.0) -- (7.55,1.50) -- cycle;

%% ================== les trois blocs ===================================
\node[bloc] (b1) at (1.85,2.0) {Convertir\\ $\eta_1$};
\node[bloc] (b2) at (4.15,2.0) {Transmettre\\ $\eta_2$};
\node[bloc] (b3) at (6.45,2.0) {Agir\\ $\eta_3$};

%% ================== puissances d'entrée et de sortie ==================
\node[font=\small, text=solideD, anchor=east] at (0.65,2.0) {$P_{\text{entrée}}$};
\node[font=\small, text=solideD, anchor=west] at (8.35,2.0) {$P_{\text{sortie}}$};

%% ================== pertes ============================================
\foreach \x in {1.85,4.15,6.45} {
  \draw[perte] (\x,1.45) -- (\x,0.85);
  \node[et, text=black!55, anchor=north] at (\x,0.85) {pertes}; }

%% ================== relation ==========================================
\node[draw=solideA, line width=1pt, fill=solideA!7, rounded corners=3pt, inner sep=5pt,
      font=\small, anchor=north] at (4.15,0.2)
     {$\eta = \dfrac{P_{\text{sortie}}}{P_{\text{entrée}}} = \eta_1 \times \eta_2 \times \eta_3$};
\node[et, anchor=north, align=center] at (4.15,-1.25)
     {exemple : $P_{\text{entrée}} = 10\,000$ W et $P_{\text{sortie}} = 6\,000$ W
      $\;\Rightarrow\;$ $\eta = \dots$};
\end{tikzpicture}
\end{document}
